\documentclass[12pt,a4paper,titlepage]{article}
\usepackage[utf8]{inputenc}
% \usepackage[francais]{babel}

\usepackage[left=1.5cm,right=2cm,top=2cm,bottom=2cm]{geometry}
\usepackage{amsmath}
\usepackage{amssymb}
\usepackage{amsthm}
\usepackage{amsfonts}
\usepackage{xcolor}
\usepackage[colorlinks]{hyperref}
\usepackage{graphicx}

% --- CSL part
\usepackage{natbib}
%\input{/project/pvs/doc/prelude}
\def\labelitemii{$\circ$}
\def\labelitemiii{$\star$}
\def\labelitemiv{$\diamond$}
\newcommand{\tcc}{{\small\small TCC}}
\newcommand{\tccs}{\tcc s}

%\renewcommand{\memo}[1]{\mbox{}\par\vspace{0.25in}\noindent\fbox{\parbox{.95\linewidth}{{\bf Memo: }#1}}\vspace{0.25in}}

\newcommand{\eg}{{\em e.g.\/},}
\newcommand{\ie}{{\em i.e.\/},}

\newcommand{\pvs}{PVS}

\newcommand{\ch}{\choice}
\newcommand{\rsv}[1]{{\rm\tt #1}}

\newcommand{\lpvstheory}[3]{\figurehead{\hozline\smaller\smaller\begin{alltt}}%
                           \figuretail{\end{alltt}\vspace{-0in}\hozline}%
                           \figurelabel{#3}\figurecap{#2}%
                           \begin{longfigure}\input{#1}\end{longfigure}}

\newcommand{\pvstheory}[3]
  {\begin{figure}[htb]\begin{boxedminipage}{\textwidth}%
   {\smaller\smaller\begin{alltt} \input{#1}\end{alltt}}\end{boxedminipage}%
   \caption{#2}\label{#3}\end{figure}}

\newcommand{\bpvstheory}[3]
  {\begin{figure}[b]\begin{boxedminipage}{\textwidth}%
   {\smaller\smaller\begin{alltt} \input{#1}\end{alltt}}\end{boxedminipage}%
   \caption{#2}\label{#3}\end{figure}}

\newcommand{\spvstheory}[1]
  {\vspace{0.1in}\par\noindent\begin{boxedminipage}{\textwidth}%
   {\smaller\smaller\begin{alltt} \input{#1}\end{alltt}}\end{boxedminipage}\vspace{0.1in}%
   }
%\newenvironment{spvstext}%
%  {\vspace{0.1in}\par\noindent\begin{boxedminipage}{\textwidth}%
%   \smaller\smaller\begin{alltt}}%
%  {\end{alltt}\end{boxedminipage}\vspace{0.1in}%
%   }


%  {\begin{boxedminipage}{\textwidth}{\smaller\smaller\begin{alltt}#1\end{alltt}}\end{boxedminipage}}
%\newenvironment{spvstheory}{\par\noindent\begin{boxedminipage}{\textwidth}\smaller\smaller\begin{alltt}}{\end{alltt}\end{boxedminipage}}

\newenvironment{pvsex}%
  {\setlength{\topsep}{0in}\smaller\begin{alltt}}%
  {\end{alltt}}

\newcommand{\pvsbnf}[2]
  {\begin{figure}[htb]\begin{boxedminipage}{\textwidth}%
   \input{#1}\end{boxedminipage}\caption{#2}\label{#1}\end{figure}}

\newcommand{\spvsbnf}[1]
  {\begin{boxedminipage}{\textwidth}\input{#1}\end{boxedminipage}}

\newcommand{\pidx}[1]{{\rm #1}} % primary index entry
\newcommand{\sidx}[1]{{\rm #1}} % secondary index entry
\newcommand{\cmdindex}[1]{\index{#1@\cmd{#1}}}
\newcommand{\icmd}[1]{\cmd{#1}\cmdindex{#1}}
\newcommand{\iecmd}[1]{\ecmd{#1}\cmdindex{#1}}
\newcommand{\buf}[1]{\texttt{#1}}
\newcommand{\ibuf}[1]{\buf{#1}\index{#1 buffer@\buf{#1} buffer}\index{buffers!\buf{#1}}}

\newenvironment{pvscmds}%
  {\par\noindent\smaller%
   \begin{tabular*}{\textwidth}{|l@{\extracolsep{\fill}}l@{\extracolsep{\fill}}l|}\hline%
     {\it Command} & {\it Aliases} & {\it Function}\\ \hline}%
  {\hline\end{tabular*}\vspace{0.1in}}

\newenvironment{pvscmdsna}%
  {\par\noindent\smaller%
   \begin{tabular*}{\textwidth}{|l@{\extracolsep{\fill}}l|}\hline%
     {\it Command} & {\it \,\,Function}\\ \hline}%
  {\hline\end{tabular*}\vspace{0.1in}}

\newcommand{\cmd}[1]{{\tt #1}}
\newcommand{\ecmd}[1]{{\tt M-x #1}}

\newcommand{\latex}{\LaTeX}                  %  LaTeX
\newcommand{\sun}{{S{\smaller\smaller UN}}}                 %  Sun
\newcommand{\sparc}{{S{\smaller\smaller PARC}}}             %  Sparc
\newcommand{\sunos}{{S{\smaller\smaller UN}OS}}             %  SunOS
\newcommand{\solaris}{{\em Solaris\/}}        %  Solaris
\newcommand{\sunview}{{S{\smaller\smaller UN}V{\smaller\smaller IEW}}} %SunView
\newcommand{\unix}{{U{\smaller\smaller NIX}}}               %  Unix
\newcommand{\lisp} {{\sc Lisp}}              %  Lisp
\newcommand{\gnu}{{Gnu Emacs}}           %  Gnu Emacs
\newcommand{\gnuemacs}{{Gnu Emacs}}      %  Gnu Emacs
\newcommand{\emacsl}{{Emacs-Lisp}}       %  Emacs Lisp
\newcommand{\shell}{{\sc Csh}}               %  C-shell

\newcommand{\update}[3]{#1\{#2\leftarrow #3\}}
\newcommand{\interp}[3]{\cal{M}\dlb {\tt #1 : #2 }\drb #3}
\newcommand{\myforall}[2]{(\forall{#1 .}\ #2)}
\newcommand{\myexists}[2]{(\exists{#1 .}\ #2)}
\newcommand{\mth}[1]{$ #1 $}
\newcommand{\labst}[2]{(\lambda{#1}.\ #2)}
\newcommand{\app}[2]{(#1\ #2)}
\newcommand{\problem}[1]{{\bf Exercise: } {\em #1}}
\newcommand{\rectype}[1]{[\# 1 \#]}
\newcommand{\recttype}[1]{{\tt [\# 1 \#]}}
\newcommand{\dlb}{\lbrack\!\lbrack}
\newcommand{\drb}{\rbrack\!\rbrack}
\newcommand{\cross}{\times}
\newcommand{\key}[1]{{\tt #1}}
\newcommand{\keyindex}[1]{\index{#1@\key{#1}}}
\newcommand{\ikey}[1]{\key{#1}\keyindex{#1}}
\newcommand{\keyword}[1]{{\smaller\texttt{#1}}}

\newenvironment{keybindings}%
  {\begin{center}\begin{tabular}{|l|l|}\hline Key & Function\\ \hline}%
  {\hline\end{tabular}\end{center}}
\def\rmif{\mbox{\bf if\ }}
\def\rmiff{\mbox{\bf \ iff \ }}
\def\rmthen{\mbox{\bf \ then }}
\def\rmelse{\mbox{\bf \ else }}
\def\rmend{\mbox{\bf end}}
\def\rmendif{\mbox{\bf \ endif}}
\def\rmotherwise{\mbox{\bf otherwise}}
\def\rmwith{\mbox{\bf \ with\ }}
\def\mapb{\char"7B\char"7B}
\def\mape{\char"7D\char"7D}
\def\setb{\char"7B}
\def\sete{\char"7D}

% ---------------------------------------------------------------------
% Macros for little PVS sessions displayed in boxes.
%
% Usage: (1) \setcounter{sessioncount}{1} resets the session counter
%
%        (2) \begin{session*}\label{thissession}
%             .
%              < lines from PVS session >
%             .
%            \end{session*}
%
%            typesets the session in a numbered box in ALLTT mode.
%
%  session instead of session* produces unnumbered boxes
%
%  Author: John Rushby
% ---------------------------------------------------------------------
\newlength{\hsbw}
\newenvironment{session}{\begin{flushleft}
 \setlength{\hsbw}{\linewidth}
 \addtolength{\hsbw}{-\arrayrulewidth}
 \addtolength{\hsbw}{-\tabcolsep}
 \begin{tabular}{@{}|c@{}|@{}}\hline 
 \begin{minipage}[b]{\hsbw}
% \begingroup\small\mbox{ }\\[-1.8\baselineskip]\begin{alltt}}{\end{alltt}\endgroup\end{minipage}\\ \hline
 \begingroup\sessionsize\vspace*{1.2ex}\begin{alltt}}{\end{alltt}\endgroup\end{minipage}\\ \hline
 \end{tabular}
 \end{flushleft}}
\newcounter{sessioncount}
\setcounter{sessioncount}{0}
\newenvironment{session*}{\begin{flushleft}
 \refstepcounter{sessioncount}
 \setlength{\hsbw}{\linewidth}
 \addtolength{\hsbw}{-\arrayrulewidth}
 \addtolength{\hsbw}{-\tabcolsep}
 \begin{tabular}{@{}|c@{}|@{}}\hline 
 \begin{minipage}[b]{\hsbw}
 \vspace*{-.5pt}
 \begin{flushright}
 \rule{0.01in}{.15in}\rule{0.3in}{0.01in}\hspace{-0.35in}
 \raisebox{0.04in}{\makebox[0.3in][c]{\footnotesize \thesessioncount}}
 \end{flushright}
 \vspace*{-.57in}
 \begingroup\small\vspace*{1.0ex}\begin{alltt}}{\end{alltt}\endgroup\end{minipage}\\ \hline 
 \end{tabular}
 \end{flushleft}}
\def\sessionsize{\footnotesize}
\def\smallsessionsize{\small}
\newenvironment{smallsession}{\begin{flushleft}
 \setlength{\hsbw}{\linewidth}
 \addtolength{\hsbw}{-\arrayrulewidth}
 \addtolength{\hsbw}{-\tabcolsep}
 \begin{tabular}{@{}|c@{}|@{}}\hline 
 \begin{minipage}[b]{\hsbw}
 \begingroup\smallsessionsize\mbox{ }\\[-1.8\baselineskip]\begin{alltt}}{\end{alltt}\endgroup\end{minipage}\\ \hline 
 \end{tabular}
 \end{flushleft}}
\newenvironment{spec}{\begin{flushleft}
 \setlength{\hsbw}{\textwidth}
 \addtolength{\hsbw}{-\arrayrulewidth}
 \addtolength{\hsbw}{-\tabcolsep}
 \begin{tabular}{@{}|c@{}|@{}}\hline 
 \begin{minipage}[b]{\hsbw}
 \begingroup\small\mbox{
}\\[-0.2\baselineskip]}{\endgroup\end{minipage}\\ \hline 
 \end{tabular}
 \end{flushleft}}
\newcommand{\memo}[1]{\mbox{}\par\vspace{0.25in}%
\setlength{\hsbw}{\linewidth}%
\addtolength{\hsbw}{-2\fboxsep}%
\addtolength{\hsbw}{-2\fboxrule}%
\noindent\fbox{\parbox{\hsbw}{{\bf Memo: }#1}}\vspace{0.25in}}
\newcommand{\nb}[1]{\mbox{}\par\vspace{0.25in}\setlength{\hsbw}{\linewidth}\addtolength{\hsbw}{-1.5ex}\noindent\fbox{\parbox{\hsbw}{{\bf Note: }#1}}\vspace{0.25in}}

%%% Local Variables: 
%%% mode: latex
%%% TeX-master: t
%%% End: 

\usepackage{boxedminipage}
\usepackage{tabularx}


\newcommand{\classnm}[1]{\texttt{\textup{#1}}}
\newcommand{\classhln}[1]{\hyperlink{#1}{\classnm{#1}}}
\newcommand{\classidx}[1]{\classnm{#1}\index{#1@\classnm{#1}}}
\newcommand{\classhtgt}[1]{\hypertarget{#1}{\classnm{#1}}\index{#1@\classnm{#1}|underline}}
\newcommand{\slotnm}[1]{{\itshape\ttfamily #1}}
\newcommand{\slotidx}[1]{\slotnm{#1}\index{#1@\slotnm{#1}}}
\newcommand{\slotidxi}[2]{\slotnm{#1}\index{#2@\slotnm{#2}}}
\newenvironment{class}[4]% name, kind, superclasses, slots
{\par\noindent\begin{boxedminipage}{\textwidth}\par\noindent\classhtgt{#1} %
$\subset$ \ifx!#3!$\emptyset$\else\classnm{#3}\fi%
\ifx!#2!\hfill[\emph{class}]\else\hfill[\emph{#2 class}]\fi%
\ifx!#4!\else\newline%
\begin{tabularx}{\linewidth}{l<{\dotfill}X}
#4
\end{tabularx}\fi\par\dotfill\par}%
{\end{boxedminipage}}
\usepackage{includes/makebnf}


\usepackage[space]{grffile}

% --- Partie code
\usepackage{listings}
\lstset{language=C,
        showstringspaces=false,
        basicstyle=\footnotesize\ttfamily,
        captionpos=b,
        stepnumber=1,
        keywordstyle=\bfseries\color{green!40!black},
        commentstyle=\itshape\color{purple!40!black},
        identifierstyle=\color{blue},
        stringstyle=\color{red}}

\newcommand{\cl}[1]{\texttt{#1}}

% --- Parti math. (app.)
\newtheorem{theorem}{Theorem}
\newtheorem{lemma}[theorem]{Lemma}
\newtheorem{proposition}[theorem]{Proposition}
\newtheorem{corollary}[theorem]{Corollary}
\newtheorem{definition}[theorem]{Definition}
\newtheorem{algorithm}[theorem]{Algorithm}
\newtheorem{remark}[theorem]{Remark}

% Mathematic abbreviations

\newcommand{\N}{\mathbb{N}}
\newcommand{\Z}{\mathbb{Z}}
\newcommand{\Q}{\mathbb{Q}}
\newcommand{\R}{\mathbb{R}}
\newcommand{\Esp}{\mathbb{E}}
%\newcommand{\E}{\mathbb{E}} % Si tu l'utilise beaucoup
\newcommand{\Prob}{\mathbb{P}}
%\newcommand{\P}{\mathbb{P}} % Si tu l'utilise beaucoup

\newcommand{\mpzt}{\texttt{mpz\_t}}
\newcommand{\mpqt}{\texttt{mpq\_t}}

\title{
\vspace{-3cm}
\normalsize
\begin{tabular}{p{15cm}}
ÉCOLE POLYTECHNIQUE \\
PROMOTION X2011 \\
Férey Gaspard
\end{tabular}
\vspace{3cm}
\begin{center}
\includegraphics[height=3cm]{pictures/logo.png}
\end{center}
\vspace{1cm}
\large
\begin{center}
RAPPORT DE STAGE D'OPTION SCIENTIFIQUE\\
\vspace{1cm}
{\Huge Titre}\\
\vspace{1cm}
NON CONFIDENTIEL
\end{center}
\vspace{3cm}
\normalsize
\begin{tabular}{p{6cm} p{10cm}}
Option :                  & INFORMATIQUE \\
Champ de l'option :       & Math-Informatique \\
Directeur de l'option :   & Olivier Bournez \\
Directeur de stage :      & Olivier Bournez \\
Dates du stage :          & 7 avril - 22 aout 2014\\
Nom et adresse de l'organisme :  & SRI International \\
                          & Computer Science Laboratory (CSL) \\
                          & 333 Ravenswood Avenue \\
                          & Menlo Park, CA 94025-3493 \\
                          & United States
\end{tabular}
}

\begin{document}

\maketitle
\tableofcontents
\newpage

\section{Introduction}


\section{SRI}

\subsection{The HACMS Project}

\subsection{PVS}

\subsection{Translating PVS}
\subsubsection{Parsing and typechecking PVS}
These two task we leave to PVS native parser and typechecker.

The parser generates objects representing the expressions of the theory.

We only convert a subset of PVS.
This subset is defined by a subset of expression objects we can translate.
The objective is, of course, to be able to translate the maximum of (if not all) PVS expression objects.

\subsubsection{Other translator}
- Common Lisp (native)
- Clean
- Yices




\section{Translating PVS Syntax}

\subsection{PVS Syntax}

We describe here the syntax of PVS and the objects system used to represent them in Lisp. Some slots of the classes are voluntarily omitted. For a full description of PVS parser representation, refer to \cite{PVS:api}.


%
% this (bnf-expr-aux.tex) is the latter half of bnf-expr.tex
%

\begin{bnf}

\production{Expr}{
\hlink{Number} \\[-0.1in] & \choice &
\hlink{Name} \\[-0.1in] & \choice &
Expr\ \hlink{Arguments} \\[-0.1in] & \choice &
Expr\ \hlink{Binop}\ Expr \\[-0.1in] & \choice &
\hlink{Unaryop}\ Expr \\[-0.1in] & \choice &
Expr \lit{`} \brc{ \hlink{Id} \choice \hlink{Number} } \\[-0.1in] & \choice &
\lit{(} \ites{Expr}{,} \lit{)} \\[-0.1in] & \choice &
\lit{(\#} \ites{\hlink{Assignment}}{,} \lit{\#)} \\[-0.1in] & \choice &
\hlink{IfExpr} \\[-0.1in] & \choice &
\lit{LET} \ites{\hlink{LetBinding}}{,} \lit{IN} Expr  \\[-0.1in] & \choice &
Expr \lit{WHERE} \ites{\hlink{LetBinding}}{,}  \\[-0.1in] & \choice &
Expr \lit{WITH} \lit{[} \ites{\hlink{Assignment}}{,} \lit{]} \\[-0.1in] & \choice &}

\production{Number}{\ite{\hlink{Digit}}}
\production{Id}{\hlink{Letter} \ite{\hlink{IdChar}}}
\production{IdChar}{\hlink{Letter} \choice \hlink{Digit}}
\production{Letter}{\lit{A} \choice \ldots \choice \lit{Z}}
\production{Digit}{\lit{0} \choice \ldots \choice \lit{9}}

\production{Arguments}
{\lit{(} \ites{\hlink{Expr}}{,} \lit{)}}

\production{IfExpr}
{\lit{IF} \hlink{Expr} \lit{THEN} \hlink{Expr} \\[-0.1in]
& & \rep{\brc{\lit{ELSIF} \hlink{Expr} \lit{THEN} \hlink{Expr}}} \lit{ELSE} \hlink{Expr} \lit{ENDIF}}

\production{Name}{
\lit{true}  \choice \lit{false} \choice \lit{number\_field\_pred} \choice \lit{real\_pred} \\[-0.1in]
& \choice & \lit{integer\_pred} \choice \lit{integer?} \choice \lit{rational\_pred} \\[-0.1in]
& \choice & \lit{floor} \choice \lit{ceiling} \choice \lit{rem} \choice \lit{ndiv} \choice \lit{even?} \choice \lit{odd?} \\[-0.1in]
& \choice & \lit{cons} \choice \lit{car} \choice \lit{cdr} \choice \lit{cons?} \choice \lit{null} \choice \lit{null?} \\[-0.1in]
& \choice & \lit{restrict} \choice \lit{length} \choice \lit{member} \choice \lit{nth} \choice \lit{append} \choice \lit{reverse}}

\production{Binop}{
\lit{=}  \choice \lit{\char'134 =} \choice \lit{OR} \choice \lit{\char'134 /} \choice \lit{AND} \choice \lit{\&} \choice \lit{/\char'134} \\[-0.1in]
& \choice & \lit{IMPLIES} \choice \lit{=>}  \choice \lit{WHEN} \choice \lit{IFF} \choice \lit{<=>} \\[-0.1in]
& \choice & \lit{+} \choice \lit{-} \choice \lit{*} \choice \lit{/} \choice \lit{<} \choice \lit{<=} \choice \lit{>} \choice \lit{>=}}

\production{Unaryop}{ \lit{NOT}  \choice \lit{-}}

\production{Assignment}
{ \ite{\hlink{AssignArg}} \brc{\lit{:=} \choice \lit{|->}} \hlink{Expr}}

\production{AssignArg}
{\lit{(} \ites{\hlink{Expr}}{,} \lit{)} \\[-0.1in]
& \choice & \lit{`} \hlink{Id} \\[-0.1in]
& \choice & \lit{`} \hlink{Number}}

\production{LetBinding}
{\brc{\hlink{LetBind} \choice \lit{(} \ites{\hlink{LetBind}}{,} \lit{)}} \lit{=} \hlink{Expr}}

\production{LetBind}
{ \hlink{Id} \opt{\lit{:} \hlink{TypeExpr}}}

\end{bnf}


%\production{AssignArgs}
%{\hlink{Id} \opt{\lit{!} \hlink{Number}}\\[-0.1in]
%& \choice & \hlink{Number}\\[-0.1in]
%& \choice & \ite{\hlink{AssignArg}}}

%\production{IdOp}
%{\hlink{Id} \ch \hlink{Opsym} \ch \hlink{Number}}
%
%\production{Opsym}
%{\hlink{Binop} \choice \hlink{Unaryop} \choice \lit{TRUE}
%\choice \lit{FALSE}}


%\production{BindingOp}
%{\lit{LAMBDA} \choice \lit{FORALL} \choice \lit{EXISTS} \choice \{ \hlink{IdOp} \lit{!} \}}

%\production{IdOps}
%{\ites{\hlink{IdOp}}{,}}

%\production{Selection}
%{\hlink{IdOp} \opt{\lit{(} \hlink{IdOps} \lit{)}} \lit{:} \hlink{Expr}}




\begin{class}{expr}{abstract}{\classhln{syntax}}
{\slotidx{type} & the type of the expression}
\end{class}

\begin{class}{name}{mixin}{\classhln{syntax}}
{\slotidx{id} & the identifier \\
 \slotidx{actuals} & a list of actual parameters \\
 \slotidx{resolutions} & singleton}
This is a mixin for names, i.e., \texttt{name-expr}s, \texttt{type-name}s, etc.
\end{class}

\begin{class}{name-expr}{}{\classhln{name} \classhln{expr}}{}
\end{class}

\begin{class}{number-expr}{}{\classhln{expr}}{\slotidx{number} & a nonnegative integer}
\end{class}

\begin{class}{tuple-expr}{}{\classhln{expr}}{\slotidx{exprs} & a list of expressions}
\end{class}

\begin{class}{application}{}{\classhln{expr}}%
{\slotidx{operator} & an expr \\
 \slotidx{argument} & an expr (maybe a tuple-expr)}
\end{class}

\begin{class}{field-application}{}{\classhln{expr}}%
{\slotidx{id} & identifier \\
 \slotidx{actuals} & a list of actuals \\
 \slotidx{argument} & the argument}
A field application is the internal representation for record extraction,
e.g., \texttt{r`a}
\end{class}

\begin{class}{record-expr}{}{\classhln{expr}}{\slotidx{assignments} & a list of assignments}
\end{class}

\begin{class}{lambda-expr}{}{\classhln{binding-expr}}{}
This is the subclass of \texttt{binding-expr} used for LAMBDA expressions.
\end{class}

\begin{class}{if-expr}{}{\classhln{application}}{}
When an application has an operator that resolves to the \texttt{if\_def} it is changed to this class.
\end{class}

\begin{class}{update-expr}{}{\classhln{expr}}
{\slotidx{expression} & an expr \\
 \slotidx{assignments} & a list of assignments}
An update expression of the form \texttt{e WITH [x := 1, y := 2]}, maps to
an \texttt{update-expr} instance, where the \texttt{expression} is
\texttt{e}, and the \texttt{assignments} slot is set to the list of
generated \texttt{assigment} instances.  Note that these are very
succinct representations, but correspondingly difficult to typecheck or to
translate to other systems (e.g., decision procedures).  See the
description of the function \cl{translate-update-to-if!} for more
 details.
\end{class}

\begin{class}{assignment}{}{\classhln{syntax}}
{\slotidx{arguments} & the list of arguments \\
 \slotidx{expression} & the value expression}
Assignments occur in both record-exprs and update-exprs.  The
\texttt{arguments} form is a list of lists.  For example, given the
assignment \texttt{`a(x, y)`1 := 0}, the \texttt{arguments} are
\texttt{((a) (x y) (1))} and the \texttt{expression} is \texttt{0}.
\end{class}



\subsection{Translator architecture}

Describe here the Lisp functions and data structures

Skeleton

Expected input

Output objects

Assertions that we (try to) maintain

Main steps :
\begin{itemize}
\item Typechecking: The PVS typechecker perform a type analysis on the PVS code to associate a PVS type to each expression. This might generates some proof obligations (TCC).
\item Lexical and syntactic analysis: The PVS parser transforms PVS code into a Lisp internal representation.
\item Translation: The translator generates a different representation from PVS expressions and functions declarations.
Typically, an expression $e$ is translated into a tuple of four elements $(t,n,i,d)$, where $t$ represents a C type used to describe the expression, $n$ is a string representing the expression, $i$ is a list of instructions supposed to be executed prior to using $n$ (initialisation of $n$) and $d$ is a list of instructions to be executed when $n$ isn't needed anymore (destruction of $n$).
\item Optimizations
\item Code generation: C code is generated.
\end{itemize}

We first define a function $T$ to translate an expression $e$.
$$ T(e) = ( \ T^t(e) \ , \ T^n(e) \ , \ T^i(e) \ , \ T^d(e) \ ) $$

\begin{figure}[h!]
\begin{eqnarray*}
T(\cl{2}) &= (& \cl{int}, "2", [], []) \\
T(\cl{4294967296}) &= (& \cl{mpz\_t}, \ ? \ , \\
&& [ \cl{mpz\_init(?);} \ | \\
&& \ \cl{mpz\_set\_str(?, "4294967296");} ],\\
&& [ \cl{mpz\_clear(?);} ]) \\
T(\cl{lambda(x:below(10)):x}) &= (& \cl{int*}, \ ? \ , \\
&& [ \cl{? = malloc(10 * sizeof(int));} \ | \\
&& \ \cl{int i;} | \\
&& \ \cl{for(i = 0; i < 10; i++) }\\
&& \ \ \ \cl{?[i] = i;}  ]\\
&& [ \cl{free(?);} ])
\end{eqnarray*}
\caption{Translation examples: number expressions}
\end{figure}

It may occur that $T^n(e) = ?$. In that case, the symbol $?$ appearing in $T^i(e)$ and $T^d(e)$ needs to be replaced by a proper variable name.\\

We then define two other operators:
\begin{itemize}
\item $R$ wich take an expression and a type and may add an extra conversion in the instructions to make sure its result has the expected type. Also the result of this function has a proper name.
\item $S$ which take an expression, a type and a name. It makes sure that the given variable (type + name) is set to a value representing the expression.
\end{itemize}




\subsection{A few translation rules}

Translation rules :




\begin{lstlisting}
number-expr "2"
(C-int, "2", [], [])

number-expr "12315468453213"
(C-mpz, nil,
        [mpz_t ~a; | mpz_t_init("12315468453213"); ],
        [mpz_clear ~a;])

application "f(e1, e2)"
(C-mpz, nil,
       [ instr(e1) | instr(e2)
                   | mpz(~a); | f(~a, e1, e2) ]
       [mpz_clear(~a);]
\end{lstlisting}



\section{Types}

\subsection{PVS Types}
A PVS theory can be typechecked using the emacs interface \cl{M-x typecheck} or with Lisp function \cl{(tc name-theory)}. This first runs the PVS parser on the code and generates CLOS objects to represent it. Then, the PVS typechecker is run on this internal representation of the theory and tries to give a type to all expressions generating TCC when needed.\\

Here we describe how PVS types are represented in Lisp. The syntax of PVS we allow


\begin{bnf}

\hproduction{type-expr}{TypeExpr}
{\hlink{Name} \\[-0.1in]
& \choice & \hlink{EnumerationType} \\[-0.1in]
& \choice & \hlink{Subtype} \\[-0.1in]
& \choice & \hlink{TypeApplication} \\[-0.1in]
& \choice & \hlink{FunctionType} \\[-0.1in]
& \choice & \hlink{TupleType} \\[-0.1in]
& \choice & \hlink{CotupleType} \\[-0.1in]
& \choice & \hlink{RecordType}}

\hproduction{enumtype}{EnumerationType}
{\lit{\setb} \hlink{IdOps} \lit{\sete}}

\hproduction{subtype}{Subtype}
{\lit{\setb} \hlink{SetBindings} \lit{|} \hlink{Expr} \lit{\sete} \\[-0.1in]
& \choice & \lit{(} \hlink{Expr} \lit{)} }

\hproduction{type-application}{TypeApplication}
{\hlink{Name}\ \hlink{Arguments}}

\hproduction{funtype}{FunctionType}
{\opt{\lit{FUNCTION} \choice \lit{ARRAY}} \\[-0.1in]
& & \lit{[} \ites{\setb \opt{\hlink{IdOp} \lit{:}} \hlink{TypeExpr} \sete}{,} \lit{->} \hlink{TypeExpr} \lit{]}}

\hproduction{tupletype}{TupleType}
{\lit{[} \ites{\setb \opt{\hlink{IdOp} \lit{:}} \hlink{TypeExpr} \sete}{,} \lit{]}}

\hproduction{cotupletype}{CotupleType}
{\lit{[} \ites{\setb \opt{\hlink{IdOp} \lit{:}} \hlink{TypeExpr} \sete}{+} \lit{]}}

\hproduction{recordtype}{RecordType}
{\lit{[\#} \ites{\hlink{FieldDecls}}{,} \lit{\#]}}

\hproduction{field-decl}{FieldDecls}
{\hlink{Ids}\ \lit{:} \hlink{TypeExpr}}
\end{bnf}


\begin{class}{type-expr}{abstract}{\classhln{syntax}}
{}
\end{class}
\begin{class}{type-name}{}{\classhln{type-expr} \classhln{name}}
{\slotidx{adt} & }
\end{class}
\begin{class}{subtype}{}{\classhln{type-expr}}
{\slotidx{supertype} & \\
 \slotidx{predicate} & }
\end{class}
\begin{class}{funtype}{}{\classhln{type-expr}}
{\slotidx{domain} & \\
 \slotidx{range} & }
\end{class}
\begin{class}{tupletype}{}{\classhln{type-expr}}
{\slotidx{types} & }
\end{class}
\begin{class}{recordtype}{}{\classhln{type-expr}}
{\slotidx{fields} & }
\end{class}


\subsection{Translating types}
PVS types:\\
boolean, number, number\_field, real, 
rational, integer, $A \rightarrow B$, restricted types
$\text{below}(10) := \{ x:\text{int} | 0 \leq x < 10 \} $)  \\
enum
datatype


This requires a type analysis to decide on the type of a PVS expression. For example the PVS \cl{int} type can be represented by the \cl{int}, \cl{unsigned long} or \cl{mpz\_t} C types. In that case, we study the range of the expression to decide which types are allowed to represent it. Then we take the context in which the expression appears to decide. For instance in
\begin{lstlisting}
incr(x:below(10)):int = x+1
\end{lstlisting}
the \cl{x} expression, result of the function \cl{incr} can always be represented by an \cl{int} or \cl{unsigned long} in C but we choose here to represent it using a \cl{mpz\_t}.

Auxiliary type system :
C-type with a flag : mutable (meaning that the expression it describes only has one pointer pointing to it.

\begin{lstlisting}
int a = 2;     a : int[mutable]
int* a = malloc( 10 * sizeof(int*) );  
\end{lstlisting}

destructive addition:

\begin{lstlisting}
d_add(*mpz_t res, mpz_t[mutable] a, long b) {
  mpz_add(a, a, b);
  (*res) = a;
}
\end{lstlisting}

Rq : \cl{d\_add} is given a mutable \mpzt, meaning that it can modify it and is responsible for freeing it.
It is also responsible for allocating memory for the result.
Here it uses the memory to assign res.



Use an auxiliary language :

\begin{lstlisting}
( expr, C-type[mutable] )
\end{lstlisting}

Conversions and copies create mutables types (at a cost) :
\cl{a[mutable]\_from\_b}

\cite{fillcomp}







C types:\\
\cite{huss2004c}
\begin{lstlisting}[caption=C types]
// integer and floating point types
[unsigned] char, int, long, double
type*  //arrays
char*  // strings
struct types  // structures with fields
enum types
short int, float, union, size_t // etc...
\end{lstlisting}










We can only translate a subset of all PVS types.
What's missing ?




\subsection{Translating PVS syntax}


We can only translate a subset of PVS syntax.
What's missing ?



\section{Difficulties and successes}


\subsection{\texttt{if} expressions}
Represented by \texttt{if-expr}

\subsection{Integer, rationnals}
In PVS, the \texttt{integer} represent the whole set $\Z$ of all relative numbers (and \texttt{rational} also describe $\Q$).
In C, we have finite types
int, long, ...

We need the GMP library which introduces the types \mpzt and \mpqt.
These types are arrays and should be used just as integer (not as pointers except they still need to be freed).


\subsection{Garbage collection }

We implement a very simple "Reference Counting Garbage Collector" as described in \cite{jonesgarbage}.\\

We maintain a hashtable of pointer counters.
Each pointer in the code is a key in the hastable to which we associate an int counter as value.\\

Pointers only occurs in arrays or struct.\\

Arrays are created in the code.\\


\begin{lstlisting}
T* a = b;
\end{lstlisting}
becomes
\begin{lstlisting}
T* a = GC( b ); // Should not happen often...
\end{lstlisting}

\begin{lstlisting}
t[0] = b;  // with b of type T*
\end{lstlisting}
becomes
\begin{lstlisting}
GC_free( t[0] );
t[0] = (T*) GC( b );
\end{lstlisting}

Examples

\begin{lstlisting}
int* f() {
  int* res;
  res = (int*) GC_alloc( 10 * sizeof(int) );
  [... init res...]
  return res; // pointer count = 1
}

void main() {
   int* a = f();  // pointer counter of a = 1
   int** b = (int**) GC_alloc( sizeof( int*) );
   GC_free( b[0] );  // useless
   b[0] = (int*) GC( a );  // pointer counter of a = 2
   printf("f(0) = %s", b[0][0]);
   GC_free(b); // frees b, pointer count of a = 1
   GC_free(a); // frees a
}
\end{lstlisting}


\subsection{Performance}

\subsection{Update expressions}
Update expressions are represented by PVS as \texttt{update-expr} objects.

$$ E := \texttt{ t with [ e1 := e2 ] } $$


Problem :
\cl{t} is an expression typed as a function. Therefore it might be represented in C as an array (if domain type is \cl{below($n$)}.
We want to know if we can update \cl{t} in place to obtain a C object representing \cl{E} or if we have to make a copy of \cl{t}.


We consider three solutions to this problem.

\subsubsection{Pointer counting}

We keep track of the number of pointer pointing to an array or a struct.

This requires to build our own C struct (heavy)
\begin{lstlisting}
struct array_int {
   int pointer_count = 1;
   int *data;
};
\end{lstlisting}

When we update the struct, if the pointer is 0, we update in place.\\

Besides every update require now to read the structure and make a test (small compared to a copy but no so small compared to a single in place update)\\

Besides, the creation / destruction gets more complicated\\

Passing argument to function :
\begin{lstlisting}
array_int f(array_int arg) {
   arg.pointer_count++;  \\ Since now f also have a pointer to the struct
   if (arg.pointer_count == 1) {
      arg.data[0] = 0;
      return arg;
   } else {
      array_int res;
      res.data = malloc( 10 * sizeof(int*) );
      copy(res, arg);  // Very long...
      res.pointer_count --; // This function is about to lose its pointer to res
      return res;
   }
}

void main() {
   array_int t;
   init(t);  // somehow...
   
   t.pointer_counter --;  // We assure we won't use the pointer "t" to the array anymore
   array_int r = f(t);
   t = null; // This way we garantee the variable "t" won't be used later in the code
   
   [...]   
}
\end{lstlisting}

This add quite some code compared to the simple :

\begin{lstlisting}
array_int f(array_int arg) {
   arg[0] = 0;
   return arg;
}

void main() {
   array_int t;
   init(t);  // somehow...
   array_int r = f(t);
   
   [...]   
}
\end{lstlisting}



\subsubsection{Using a different data structure}

PVS uses arrays in a very particular way, we might then represent them with an other structure than just only a C array.
For example :

\begin{lstlisting}
struct r_list_int {
   int key;
   int value;
   r_list_int tl;
};
struct array_int {
   int *data;
   r_list_int replacement_list;
};
\end{lstlisting}

Each structure represent the array \cl{data} with the modifications contained in the linked list \cl{r\_list\_int}

Problems : Just as the previous solution :
- add some extra code
- add some extra computation (runtime tests, reading the replacement list)
- require to create as many structures and associated functions as there are range types fo the manipulated arrays
- Very dependent on the GC


\subsubsection{Flow analysis on the PVS code}

\cite{pavol}



\subsubsection{Analysis of the C code}

Trying to avoid copying arrays by analyzing the C code generated.

2 different functions (destructive and non destructive)


Algorithm :

Always have a "non destructive" version of any function. A "cautious" version that never modify the arguments in place and always make copies when necessary (when a "mutable" version of an array is necessary (for instance updates)).

In destructive versions of all functions :
Flag all array arguments to "mutable".
Then for each of these arguments :
  - If it never occurs destructively, then remove flag
          (function just observe the arg)
  - If it occurs destructively, it can never occur at all AFTER.
  -> Need to define the order of evaluation of expression
      (easy rules on simple expressions)
  -> Need to be able to detect occurences of a name-expr
  -> Otherwise, unflag the arg

What is a destructive occurence :
$$ E :=  \cl{f(  t with [ e1 := e2 ] , t(0) )} $$
order of eval :
e1 and e2  (t can occur non destr)
t          ( expression of an update : destr)
t(0)       ( occurence of t (even non destr))

\cl{f(x:Arr):int = g( h(t), t)}  is destructively translated to \\
\begin{lstlisting}[numbers=left,caption=Example]
int f_d(int* t) {   // t has type ! since this is destructive f
  int* arg1 = h(t); // h can't be called destructively because
                    // even though t is !, it appears later (line 4)
  int* arg2 = t;    // t is ! and never appears later => arg2 is !
  return g( arg1, arg2); // arg2 is ! but g can only be called
}                        // destructively if arg1 is
\end{lstlisting}

if g has type \cl{[Array! -> ?]} then t can't be destructive\\

if g has type \cl{[Array -> ?]} then t can be destructive

First algorithm:


Need multiple passes as the flags disappear



\section{Conclusion}

\subsection{What's left to be done ?}

Use a C structure to represent a closure
\begin{lstlisting}
struct r_list_int {
   int (*body)(void* env, void* args);
   void* env;
   void* args;
};
\end{lstlisting}



\subsection{My stay at SRI}

Discovering PVS :
Translating Coq proofs to PVS
PVS library for basic linear algebra

Robin project, HACMS \\
Contest week-end 14-15 June \\
Summer School \\
Parsing Lisp code -> generate HTML architecture file\\
Correcting translator PVS to SMT-LIB
\cite{pavol}



\bibliographystyle{plain}
\bibliography{report}


\end{document}
